% Manual de usuario de Pymetory. Capturas tomadas de app.pymetory.com con los datos de demostración (25-sep-2026).
\section*{Manual de usuario} \label{sec:ManualesdeUsuario}
\addcontentsline{toc}{section}{Manual de usuario}

Este manual describe el uso de Pymetory en \url{https://app.pymetory.com}. Las capturas muestran los datos de demostración del caso de estudio (panadería). Las funciones marcadas como \textit{solo administrador} no están disponibles para el rol operario.

\newcommand{\capturaManual}[2]{%
\begin{figure}[H]\centering\includegraphics[width=0.92\textwidth]{Anexos/Imagenes/manual/#1}\caption{#2}\end{figure}}

\subsection*{Ingreso al sistema}
Ingrese el correo electrónico y la contraseña asignados por el administrador y pulse \textbf{Entrar}. No existe registro abierto de administradores: las cuentas con ese rol las crea otro administrador.
\capturaManual{01-login.png}{Pantalla de inicio de sesión.}

\subsection*{Tablero}
Resume el estado del inventario: número de insumos, lotes en bodega, lotes críticos por vencimiento y valorización. Debajo aparecen los últimos movimientos del Kardex y la ocupación de cada bodega. Los accesos rápidos permiten registrar un insumo, escanear un código, generar un reporte o hacer un consumo rápido.
\capturaManual{02-tablero.png}{Tablero de control.}

\subsection*{Inventario}
Muestra los lotes activos agrupados por bodega, con su cantidad, fecha de vencimiento y estado. Los lotes dentro del umbral FEFO se marcan como críticos. Desde aquí se registran ingresos (\textbf{Nuevo ingreso}) y consumos; al consumir, el sistema descuenta primero del lote que vence primero.
\capturaManual{03-inventario.png}{Inventario por lotes.}

\subsection*{Búsqueda}
Permite encontrar materiales y lotes por nombre, código o etiqueta.
\capturaManual{04-buscar.png}{Búsqueda de materiales y lotes.}

\subsection*{Asistente de inventario}
Escriba una pregunta en español, por ejemplo ``¿Qué lotes vencen pronto?'' o ``¿Cuánta harina de trigo hay?''. El asistente consulta la base de datos antes de responder; si el material no está registrado, lo indica. Las conversaciones quedan en el historial de la izquierda.
\capturaManual{05-asistente.png}{Asistente con una consulta de vencimientos.}

\subsection*{Reportes (solo administrador)}
Genera reportes de inventario, movimientos, FEFO, consumo y valorización, filtrables por bodega, material y rango de fechas, en PDF o CSV.
\capturaManual{06-reportes.png}{Módulo de reportes.}

\subsection*{Log maestro (Kardex)}
Lista todos los movimientos con tipo, material, lote, responsable, cantidad y fecha. Los movimientos no se pueden editar ni eliminar; una corrección se registra como un ajuste nuevo con su justificación.
\capturaManual{07-kardex.png}{Log maestro de movimientos.}

\subsection*{Etiquetas y escáner}
El módulo de etiquetas genera códigos QR y de barras por lote para imprimir. El escáner lee esas etiquetas con la cámara del dispositivo (o con ingreso manual del código) y registra la entrada o salida correspondiente en el Kardex.
\capturaManual{08-etiquetas.png}{Generación de etiquetas.}
\capturaManual{09-escaner.png}{Escáner de códigos QR.}

\subsection*{Transferencias}
Traslada cantidad de un lote entre bodegas. La operación registra la salida en la bodega de origen y la entrada en la de destino.
\capturaManual{10-transferencias.png}{Transferencias entre bodegas.}

\subsection*{Órdenes de compra}
Registra órdenes a proveedores. Al recibir una orden, el sistema crea el lote recibido y su entrada en el Kardex.
\capturaManual{11-ordenes.png}{Órdenes de compra.}

\subsection*{Alertas}
Reúne las alertas de stock bajo y de lotes próximos a vencer, que el sistema revisa cada hora.
\capturaManual{12-alertas.png}{Centro de alertas.}

\subsection*{Ajustes (solo administrador)}
Configura el umbral FEFO, las notificaciones, el proveedor y modelo de lenguaje del asistente, las claves de API y los usuarios con sus roles.
\capturaManual{13-ajustes.png}{Ajustes del sistema.}
